\documentclass[UTF8,fontset=none,a4paper,10pt]{ctexart}

\setmainfont{Times New Roman}
\setCJKmainfont{Songti SC}
\setCJKsansfont{Heiti SC}
\setCJKmonofont{Hiragino Sans GB}
\setmonofont{Menlo}

\usepackage[a4paper,left=19mm,right=19mm,top=18mm,bottom=19mm,headheight=14pt]{geometry}
\usepackage{amsmath,amssymb,bm,mathtools}
\usepackage{booktabs,tabularx,array,longtable}
\usepackage{xcolor}
\usepackage{enumitem}
\usepackage{listings}
\usepackage{titlesec}
\usepackage{fancyhdr}
\usepackage{microtype}
\usepackage{hyperref}
\usepackage{bookmark}

\definecolor{navy}{HTML}{17365D}
\definecolor{blue}{HTML}{2459A6}
\definecolor{lightblue}{HTML}{EAF1FB}
\definecolor{linegray}{HTML}{D7DFEA}
\definecolor{codebg}{HTML}{F4F6F8}
\definecolor{codegreen}{HTML}{26734D}

\hypersetup{
  colorlinks=true,
  linkcolor=blue,
  citecolor=blue,
  urlcolor=blue,
  bookmarksnumbered=true,
  pdfauthor={Jiguo Li},
  pdftitle={Transformer 位置编码的前世今生},
  pdfsubject={从绝对位置、相对位置到 RoPE 及长上下文扩展}
}

\setlength{\parindent}{2em}
\setlength{\parskip}{2.5pt}
\setlength{\abovedisplayskip}{6pt}
\setlength{\belowdisplayskip}{6pt}
\setlength{\abovedisplayshortskip}{4pt}
\setlength{\belowdisplayshortskip}{4pt}
\setlength{\jot}{3pt}
\linespread{1.13}
\emergencystretch=2em

\titleformat{\section}{\Large\bfseries\color{navy}}{\thesection}{0.7em}{}
\titleformat{\subsection}{\large\bfseries\color{navy}}{\thesubsection}{0.65em}{}
\titleformat{\subsubsection}{\normalsize\bfseries\color{navy}}{\thesubsubsection}{0.55em}{}
\titlespacing*{\section}{0pt}{14pt}{5pt}
\titlespacing*{\subsection}{0pt}{10pt}{3pt}
\titlespacing*{\subsubsection}{0pt}{7pt}{2pt}

\setlist[itemize]{leftmargin=1.8em,itemsep=1pt,topsep=2pt,parsep=0pt}
\setlist[enumerate]{leftmargin=2em,itemsep=1pt,topsep=2pt,parsep=0pt}

\pagestyle{fancy}
\fancyhf{}
\fancyhead[L]{\small\color{gray}Transformer 位置编码的前世今生}
\fancyhead[R]{\small\color{gray}\leftmark}
\fancyfoot[C]{\small\thepage}
\renewcommand{\headrulewidth}{0.3pt}

\lstset{
  language=Python,
  basicstyle=\ttfamily\small,
  keywordstyle=\color{blue}\bfseries,
  commentstyle=\color{codegreen},
  stringstyle=\color{purple},
  backgroundcolor=\color{codebg},
  frame=single,
  rulecolor=\color{linegray},
  breaklines=true,
  showstringspaces=false,
  columns=fullflexible,
  aboveskip=5pt,
  belowskip=5pt,
  xleftmargin=0.5em,
  framexleftmargin=0.4em
}

\newcommand{\term}[1]{\textbf{#1}}
\newcommand{\R}{\mathbb{R}}
\newcommand{\softmax}{\operatorname{softmax}}
\newcommand{\Attn}{\operatorname{Attn}}
\newcommand{\noteBox}[1]{%
  \begin{center}
  \fcolorbox{blue}{lightblue}{\parbox{0.93\linewidth}{#1}}
  \end{center}}

\begin{document}

\begin{center}
  {\zihao{1}\bfseries\color{navy} Transformer 位置编码的前世今生}\par
  \vspace{4pt}
  {\zihao{3}\color{navy} 从绝对位置、相对位置到旋转位置编码 RoPE}\par
  \vspace{10pt}
  {\normalsize\textbf{Jiguo Li\footnote{本文在Codex协助下完成}}}\par
  \vspace{2pt}
  {\small\href{mailto:jiguolee@gmail.com}{jiguolee@gmail.com}}
\end{center}
\vspace{5pt}
\hrule

\begin{abstract}
Self-Attention 能有效建模 token 之间的内容相关性，但其计算本身不包含序列顺序。位置编码的作用，是把绝对坐标、相对距离或旋转相位注入注意力，使模型能够区分词序、表达局部结构，并在更长上下文中保持可用的位置信号。

本文从 Self-Attention 的置换等变性出发，依次推导正弦--余弦位置编码、Shaw 相对位置表示、Transformer-XL、T5 Relative Position Bias 与 RoPE，比较它们注入位置的环节、计算代价、长度外推能力及对 KV Cache 的影响。在此基础上，进一步梳理 Position Interpolation、RoPE scaling law、NTK-aware、Dynamic NTK、NTK-by-parts、YaRN、LongRoPE 和 LongRoPE2，说明这些方法如何调整频率、位置尺度与注意力温度，以及训练长度、目标上下文长度和微调数据之间的约束。全文同时给出公式推导、直观解释、实现要点与一手文献，便于在模型设计和长上下文扩展中据此选择方案。
\end{abstract}

\noindent\textbf{关键词：}Transformer；位置编码；相对位置编码；RoPE；长上下文；长度外推

\noteBox{\textbf{阅读提要：}绝对位置编码为 token 指定坐标；相对位置编码描述 token 对之间的距离；RoPE 则把位置写入 Query/Key 的旋转相位。}

\clearpage
\begingroup
\small
\setlength{\parskip}{0pt}
\linespread{1.02}\selectfont
\setcounter{tocdepth}{2}
\tableofcontents
\endgroup
\clearpage

\section{为什么 Transformer 必须显式注入位置}

\subsection{RNN、CNN 与 Transformer 的差异}

RNN 按时间递推：
\begin{equation}
  h_t=f(x_t,h_{t-1}),
  \label{eq:rnn}
\end{equation}
计算路径本身就是 $x_1\rightarrow x_2\rightarrow\cdots\rightarrow x_n$，顺序天然存在于状态传递中。CNN 的卷积核先聚合局部邻域，因此也带有局部结构先验。

Transformer 去掉递归和卷积，让所有 token 并行交互。这是高吞吐的来源，也意味着模型结构本身不再知道 token 的先后顺序。原始 Transformer 因此必须额外注入位置信息\cite{vaswani2017}。

\subsection{不带位置的 Self-Attention 看到了什么}

给定输入矩阵：
\begin{equation}
  X=[x_1,x_2,\ldots,x_n]^\top\in\R^{n\times d},
  \label{eq:input}
\end{equation}
标准 Self-Attention 计算：
\begin{align}
  Q&=XW_Q,\quad K=XW_K,\quad V=XW_V, \label{eq:qkv}\\
  \Attn(X)&=\softmax\!\left(\frac{QK^\top}{\sqrt{d_h}}\right)V. \label{eq:attention}
\end{align}
第 $i$ 个 Query 对第 $j$ 个 Key 的分数为：
\begin{equation}
  s_{ij}=\frac{q_i^\top k_j}{\sqrt{d_h}}.
  \label{eq:score}
\end{equation}
该分数使用了内容向量，却没有直接使用位置 $i,j$。设 $P$ 为任意置换矩阵，则：
\begin{equation}
  \Attn(PX)=P\Attn(X).
  \label{eq:permutation}
\end{equation}
因此，不含位置机制的 Self-Attention 是\term{置换等变}的：打乱输入后，输出只按相同方式打乱。位置可以在输入表示、注意力分数或 Query/Key 的几何变换中注入，分别对应绝对位置、相对位置和 RoPE 三类方案。

\section{技术演进概览}

\begin{table}[htbp]
\centering
\small
\begin{tabularx}{\linewidth}{@{}p{1.2cm}p{3.5cm}X@{}}
\toprule
时间 & 代表工作 & 核心位置机制 \\
\midrule
2017 & Transformer & 固定正弦--余弦绝对位置编码 \\
2018 & BERT & 可学习绝对位置 embedding \\
2018 & Shaw et al. & 在注意力中加入相对距离表示 \\
2019 & Transformer-XL & 相对位置分解，支持片段级记忆 \\
2020 & T5 & 每个注意力头学习分桶距离 bias \\
2021 & RoFormer / RoPE & 按位置旋转 Query 和 Key \\
2021 & ALiBi & 对 logits 加线性距离惩罚 \\
2023 & PI / YaRN & RoPE 插值、分频缩放与注意力校准 \\
2024--2025 & LongRoPE / LongRoPE2 & 搜索逐维缩放因子并混合长短上下文训练 \\
\bottomrule
\end{tabularx}
\caption{Transformer 位置机制的代表性演进。}
\label{tab:timeline}
\end{table}

这些方法不是简单的新旧替代。固定长度 encoder、encoder-decoder 与自回归 LLM 的任务结构、推理方式和上下文需求不同，因此多种方案至今并存。

\section{第一代：绝对位置编码}

绝对位置编码为位置 $i$ 提供向量 $p_i$：
\begin{equation}
  h_i^{(0)}=e(x_i)+p_i.
  \label{eq:absolute}
\end{equation}
它相当于给每个 token 贴上“第 137 位”这样的全局坐标。

\subsection{可学习绝对位置 embedding}

最直接的实现是维护参数表：
\begin{equation}
  P\in\R^{L_{\max}\times d}.
  \label{eq:learned-table}
\end{equation}
第 $i$ 行就是位置 $i$ 的向量 $p_i$。BERT 将 token、segment 与 position embedding 相加作为输入\cite{devlin2019}。它实现简单、可直接拟合任务内位置模式，但参数表绑定最大长度；未训练位置没有可靠表示；相邻位置差不被结构性约束；内容与位置在第一层前即混合：
\begin{equation}
  q_i=(e(x_i)+p_i)W_Q=e(x_i)W_Q+p_iW_Q.
  \label{eq:absolute-proj}
\end{equation}

\subsection{正弦--余弦位置编码}

原始 Transformer 用多组正弦和余弦函数生成位置向量\cite{vaswani2017}：
\begin{align}
  PE(pos,2i)&=\sin\!\left(\frac{pos}{10000^{2i/d_{\text{model}}}}\right), \label{eq:sin-pe}\\
  PE(pos,2i+1)&=\cos\!\left(\frac{pos}{10000^{2i/d_{\text{model}}}}\right). \label{eq:cos-pe}
\end{align}
定义第 $i$ 组角频率：
\begin{equation}
  \theta_i=10000^{-2i/d_{\text{model}}},
  \label{eq:sin-theta}
\end{equation}
则一组二维表示为：
\begin{equation}
  PE_i(pos)=
  \begin{bmatrix}
    \sin(pos\theta_i)\\
    \cos(pos\theta_i)
  \end{bmatrix}.
  \label{eq:sin-pair}
\end{equation}
完整位置向量由多组不同频率拼接：
\begin{equation}
  PE(pos)=[\sin(pos\theta_0),\cos(pos\theta_0),\ldots,
  \sin(pos\theta_{d/2-1}),\cos(pos\theta_{d/2-1})].
  \label{eq:sin-full}
\end{equation}

\subsubsection{“多只钟表”的直观解释}

每对维度可看成单位圆上的一根指针，其相位与波长分别为：
\begin{align}
  \phi_i(pos)&=pos\theta_i, \label{eq:phase}\\
  \lambda_i&=\frac{2\pi}{\theta_i}=2\pi\cdot10000^{2i/d_{\text{model}}}. \label{eq:wavelength}
\end{align}
大的 $\theta_i$ 对应“快钟”，能区分相邻位置；小的 $\theta_i$ 对应“慢钟”，描述较长尺度。单只钟周期重复，多只不同转速的钟联合后能在更大范围区分位置。常数 $10000$ 是频率基数，不是数学上唯一正确的选择。

当 $d_{\text{model}}=8$ 时，四组频率为：
\begin{equation}
  \theta_0=1,\qquad \theta_1=0.1,\qquad
  \theta_2=0.01,\qquad \theta_3=0.001.
  \label{eq:d8-frequency}
\end{equation}

\subsubsection{为什么必须同时使用 sin 和 cos}

由和角公式：
\begin{equation}
  \sin((pos+\Delta)\theta)=
  \sin(pos\theta)\cos(\Delta\theta)+
  \cos(pos\theta)\sin(\Delta\theta).
  \label{eq:addition}
\end{equation}
只保存 sin 无法通过固定线性变换得到平移结果；成对保存后有：
\begin{equation}
  PE_i(pos+\Delta)=
  \begin{bmatrix}
    \cos(\Delta\theta_i)&\sin(\Delta\theta_i)\\
    -\sin(\Delta\theta_i)&\cos(\Delta\theta_i)
  \end{bmatrix}PE_i(pos).
  \label{eq:sin-shift}
\end{equation}
矩阵只依赖位移 $\Delta$。同一频率下两个位置编码的点积为：
\begin{equation}
  PE_i(m)^\top PE_i(n)=\cos((m-n)\theta_i),
  \label{eq:sin-dot}
\end{equation}
完整编码满足：
\begin{equation}
  PE(m)^\top PE(n)=\sum_i\cos((m-n)\theta_i).
  \label{eq:sin-full-dot}
\end{equation}
所以正弦编码形式上是绝对位置表示，几何关系中却包含相对位移。原始 Transformer 实际使用式~\eqref{eq:absolute}，并未直接以式~\eqref{eq:sin-full-dot} 作为注意力分数，模型仍需训练才能利用这种结构。

\noteBox{\textbf{需要注意：}位置函数可以计算到任意 $pos$，不等于整个模型能可靠理解任意长上下文。训练长度仍约束注意力分布、状态表示与优化结果。}

\section{第二代：相对位置编码}

语言和代码中的许多关系更关心 $j-i$，而不是分别知道 $i$ 与 $j$。例如“左侧最近的定义”“距调用点最近的声明”“右侧两个 token 的参数”。

\subsection{Shaw：把相对距离加入 Key 和 Value}

Shaw 等人将 token 对的相对距离直接加入 Self-Attention\cite{shaw2018}。先裁剪距离：
\begin{equation}
  r_{ij}=\operatorname{clip}(j-i,-K,K),
  \label{eq:clip}
\end{equation}
为每个距离学习相对 Key 向量 $a^K_{r_{ij}}$：
\begin{equation}
  s_{ij}=\frac{q_i^\top(k_j+a^K_{r_{ij}})}{\sqrt{d_h}}
  =\frac{q_i^\top k_j+q_i^\top a^K_{r_{ij}}}{\sqrt{d_h}}.
  \label{eq:shaw-score}
\end{equation}
第一项表示内容匹配，第二项表示当前 Query 对某个方向与距离的偏好。Value 也可加入相对表示：
\begin{equation}
  z_i=\sum_j\alpha_{ij}(v_j+a^V_{r_{ij}}).
  \label{eq:shaw-value}
\end{equation}

\subsection{Transformer-XL：跨片段记忆中的相对位置}

Transformer-XL 允许当前片段复用上一片段隐藏状态\cite{dai2019}。若缓存携带旧片段的绝对坐标，复用到新片段会产生歧义；相对位置只依赖当前 Query 与历史 Key 的真实距离。其未归一化注意力可分解为：
\begin{equation}
  \begin{aligned}
  A_{ij}={}&q_i^\top k_j+q_i^\top W_{k,R}R_{i-j}\\
           &+u^\top k_j+v^\top W_{k,R}R_{i-j},
  \end{aligned}
  \label{eq:txl}
\end{equation}
分别对应内容--内容、内容--位置、全局内容偏置和全局位置偏置。

\subsection{T5：将相对位置压缩为标量 bias}

T5 为每个注意力头和距离桶学习一个标量\cite{raffel2020}：
\begin{equation}
  s_{ij}^{(h)}=
  \frac{q_i^{(h)\top}k_j^{(h)}}{\sqrt{d_h}}
  +b_{h,\operatorname{bucket}(j-i)}.
  \label{eq:t5-bias}
\end{equation}
距离较小时使用细粒度桶，距离较大时近似对数分桶。不同 head 可学习不同距离先验，形成“内容相似度 $+$ 距离偏好”的清晰分解。

传统相对位置方法在建模上直观，但通常要处理 $n\times n$ 个 token 对。相对信息可能进入 Key、Value、logits 或多个交叉项，增加算子融合、增量解码与高效注意力实现的复杂度。RoPE 的关键价值，就是在显式表达相对位移的同时保留标准 $QK^\top$ 形式。

\section{第三代：旋转位置编码 RoPE}

RoPE 不把位置向量加到输入，也不为 token 对显式构造相对向量，而是根据绝对位置旋转 Query 和 Key\cite{su2021}。

\subsection{从二维旋转推导核心公式}

二维旋转矩阵为：
\begin{equation}
  R(\phi)=
  \begin{bmatrix}
    \cos\phi&-\sin\phi\\
    \sin\phi&\cos\phi
  \end{bmatrix}.
  \label{eq:rotation}
\end{equation}
位置 $m$ 的旋转角为 $m\theta$：
\begin{equation}
  q'_m=R(m\theta)q_m,\qquad k'_n=R(n\theta)k_n.
  \label{eq:rope-rotate}
\end{equation}
旋转后的点积为：
\begin{equation}
  \begin{aligned}
  q_m'^\top k_n'
  &=q_m^\top R(m\theta)^\top R(n\theta)k_n\\
  &=q_m^\top R((n-m)\theta)k_n.
  \end{aligned}
  \label{eq:rope-core}
\end{equation}
这里使用了：
\begin{equation}
  R(m\theta)^\top=R(-m\theta),\qquad
  R(-m\theta)R(n\theta)=R((n-m)\theta).
  \label{eq:rotation-group}
\end{equation}
单个向量 $q'_m$ 依赖绝对位置 $m$；Query-Key 交互只通过 $n-m$ 依赖相对位置。因此，\term{RoPE 在单个向量上编码绝对位置，在点积中呈现相对位置}。

\subsection{高维 RoPE 与复数形式}

将 head dimension 两两组成二维子空间，第 $r$ 对频率为：
\begin{equation}
  \theta_r=b^{-2r/d_h},\qquad r=0,1,\ldots,d_h/2-1,
  \label{eq:rope-frequency}
\end{equation}
经典基数 $b=10000$。整体旋转矩阵为：
\begin{equation}
  R_m=\operatorname{diag}\left(
  R(m\theta_0),R(m\theta_1),\ldots,R(m\theta_{d_h/2-1})
  \right).
  \label{eq:rope-block}
\end{equation}
若把二维向量写成复数 $z_r=x_{2r}+\mathrm{i}x_{2r+1}$，RoPE 等价于：
\begin{equation}
  z'_r=z_r e^{\mathrm{i}m\theta_r}.
  \label{eq:rope-complex}
\end{equation}
Query 与 Key 的相位因子相乘后为：
\begin{equation}
  e^{-\mathrm{i}m\theta_r}e^{\mathrm{i}n\theta_r}
  =e^{\mathrm{i}(n-m)\theta_r},
  \label{eq:rope-relative-phase}
\end{equation}
即绝对相位相减后只剩相对相位。

\subsection{工程实现与配对约定}

以下实现使用相邻维度配对。某些框架采用前半维与后半维配对；两者只是排列约定不同，但权重、cos/sin cache 与旋转函数必须一致。

\begin{lstlisting}
def apply_rope(x, cos, sin):
    x_even = x[..., 0::2]
    x_odd  = x[..., 1::2]
    y_even = x_even * cos - x_odd * sin
    y_odd  = x_even * sin + x_odd * cos

    y = torch.empty_like(x)
    y[..., 0::2] = y_even
    y[..., 1::2] = y_odd
    return y
\end{lstlisting}

Value 通常不旋转，因为位置主要影响“关注谁”的权重；聚合对象仍是内容 Value。

\subsection{RoPE 与 KV Cache}

固定 RoPE 配置下，历史 Key 通常以旋转后的形式缓存：
\begin{equation}
  k'_i=R_i k_i.
  \label{eq:cached-key}
\end{equation}
生成位置 $t$ 时只需计算：
\begin{equation}
  q'_t=R_tq_t,\qquad k'_t=R_tk_t.
  \label{eq:decode-rope}
\end{equation}
历史 Key 无需重算。工程上最易出错的是 position offset：prefill、decode、padding、sequence packing、prefix cache 与 sliding window 必须对同一 token 使用一致的逻辑 position id。

\subsection{RoPE 是否保证距离越远注意力越弱}

单个频率的相对项包含：
\begin{equation}
  \cos((n-m)\theta_r),\qquad \sin((n-m)\theta_r),
  \label{eq:rope-oscillation}
\end{equation}
它们随距离周期振荡，并非单调衰减。RoFormer 讨论了多频率汇总后的远距离衰减倾向\cite{su2021}，但不能对任意固定 Query/Key、任意 head 和任意距离推出“越远分数必然越小”。

\section{RoPE 长上下文扩展：从统一缩放到逐维搜索}

\subsection{问题根源：外推引入未训练相位}

设原始训练长度为 $L_0$，目标长度为 $L_1=sL_0$，其中 $s>1$。第 $r$ 个二维子空间的周期为：
\begin{equation}
  T_r=\frac{2\pi}{\theta_r}=2\pi b^{2r/d_h}.
  \label{eq:rope-period}
\end{equation}
高频维度在 $L_0$ 内经历许多周期，局部模式训练充分；低频维度可能连一个完整周期都未覆盖。直接推理到 $L_1$ 会让部分维度进入未见过的相位区间，造成 RoPE OOD。LongRoPE2 将理论临界维度定义为周期是否超过原训练长度的分界\cite{longrope2}：
\begin{equation}
  r_{\mathrm{crit}}=\min\{r:T_r\ge L_0\}.
  \label{eq:critical-dim}
\end{equation}
扩展方法真正需要处理的是各频率的缩放方式，以及模型权重对新相位分布的适应，而不仅是增大 position id 的取值范围。

\subsection{RoPE scaling law：基数、训练长度与外推范围}

Liu 等人从周期覆盖的角度研究了 RoPE 的长度外推，并提出 RoPE-based extrapolation scaling laws\cite{liu2024scalingrope}。这里的“scaling law”描述的是旋转基数、训练长度与外推范围之间的关系，与参数量、数据量和计算量之间的经典模型 scaling law 含义不同。

对原始基数 $10000$，第 $n$ 组频率的周期为：
\begin{equation}
  T_n=\frac{2\pi}{\theta_n}=2\pi\cdot10000^{2n/d_h}.
  \label{eq:scaling-period}
\end{equation}
若预训练长度为 $T_{\mathrm{train}}$，论文把能够在训练区间内至少覆盖完整周期的维度上限写为：
\begin{equation}
  d_{\mathrm{extra}}
  =2\left\lceil
  \frac{d_h}{2}\log_{10000}\!\left(\frac{T_{\mathrm{train}}}{2\pi}\right)
  \right\rceil.
  \label{eq:critical-dimension-scaling-law}
\end{equation}
超过 $d_{\mathrm{extra}}$ 的低频维度在训练中没有观察到完整周期，超出训练长度后更容易出现相位分布偏移。以 LLaMA 2 的 $d_h=128$、$T_{\mathrm{train}}=4096$ 为例，论文计算得到 $d_{\mathrm{extra}}=92$；其余 36 个维度被视为外推时更容易失稳的部分\cite{liu2024scalingrope}。

当微调阶段把旋转基数改为 $\beta>10000$，该工作用临界维度的新周期估计外推上限：
\begin{equation}
  T_{\mathrm{extra}}
  =2\pi\,\beta^{d_{\mathrm{extra}}/d_h}.
  \label{eq:scaling-law-upper-bound}
\end{equation}
反过来，给定期望长度 $\widetilde T_{\mathrm{extra}}$，可由论文给出的关系估计所需的临界基数：
\begin{equation}
  \beta_0
  =10000^{\log_{T_{\mathrm{train}}/(2\pi)}
  \left(\widetilde T_{\mathrm{extra}}/(2\pi)\right)}.
  \label{eq:scaling-law-critical-base}
\end{equation}

论文还观察到，在固定微调长度下，将基数调小也可能改善外推，因为更多维度能在训练区间内经历更充分的三角函数变化。对应的三个相位覆盖点为：
\begin{equation}
  \beta_1=\frac{2T_{\mathrm{train}}}{\pi},\qquad
  \beta_2=\frac{T_{\mathrm{train}}}{\pi},\qquad
  \beta_3=\frac{T_{\mathrm{train}}}{2\pi}.
  \label{eq:scaling-law-small-base}
\end{equation}
这组结果揭示了 RoPE 外推与周期覆盖之间的联系，但不宜把式~\eqref{eq:scaling-law-upper-bound} 当成跨模型的硬性上下文上限。论文主要在 LLaMA 2 上验证，实际有效长度还受到语料中的长距离依赖分布、微调数据、attention head 行为和评测任务影响。

\subsection{Position Interpolation：统一压缩位置索引}

PI 将目标范围中的位置压回原训练区间\cite{chen2023pi}：
\begin{equation}
  m'=\frac{m}{s}=m\frac{L_0}{L_1}.
  \label{eq:pi-position}
\end{equation}
等价地，对所有频率统一缩放：
\begin{equation}
  \theta'_r=\frac{\theta_r}{s}.
  \label{eq:pi-frequency}
\end{equation}
其优点是所有新位置都落回已训练相位范围，结构改动极小；缺点是高频也被压缩，相邻 token 的相位差从 $\theta_r$ 变成 $\theta_r/s$，局部分辨率下降。PI 论文从 2K 扩至 32K，仅需短程微调即可适应，但也明确指出原长度内会有一定性能损失\cite{chen2023pi}。

\subsection{NTK-aware：通过修改 RoPE base 非均匀缩放频率}

NTK-aware 的直觉是：不要像 PI 一样等比例压缩所有频率，而是调整频率基数 $b$，使高频尽量保持、低频承担更多扩展。常见形式把新基数设为：
\begin{equation}
  b'=b\,s^{d_h/(d_h-2)},
  \label{eq:ntk-base}
\end{equation}
从而：
\begin{equation}
  \theta'_r=(b')^{-2r/d_h}
  =\theta_r\,s^{-2r/(d_h-2)}.
  \label{eq:ntk-frequency}
\end{equation}
当 $r=0$ 时最高频保持不变；随 $r$ 增大，低频被更强地缩放。它在局部精度和长程覆盖之间比统一 PI 更平滑。需要注意，NTK-aware 最初来自开源社区经验，后来由 YaRN 论文系统整理，而不是一开始就有独立同行评议论文\cite{peng2024yarn}。

\subsection{Dynamic NTK：按当前序列长度动态选择缩放}

固定目标长度 $L_1$ 的缩放会使短输入也承受位置压缩。Dynamic Scaling 根据当前序列长度 $\ell$ 更新倍率\cite{peng2024yarn}：
\begin{equation}
  s(\ell)=\max\left(1,\frac{\ell}{L_0}\right),
  \label{eq:dynamic-scale}
\end{equation}
再将 $s(\ell)$ 代入式~\eqref{eq:ntk-base}。短序列保持原 RoPE，超过 $L_0$ 后逐步扩展。

这一做法与旋转后 KV Cache 存在冲突：当 $s(\ell)$ 变化时，同一个历史 token 的旋转角也变化。若缓存的是已经旋转的 Key，旧缓存便与新倍率不一致。严格实现要么缓存旋转前 Key 并在倍率变化后重新旋转，要么在一次生成会话内固定倍率；YaRN 论文也明确提醒了这一点\cite{peng2024yarn}。

\subsection{NTK-by-parts：按频率区间选择插值策略}

YaRN 用原训练长度内的旋转圈数衡量第 $r$ 个频率的训练覆盖：
\begin{equation}
  \rho(r)=\frac{L_0\theta_r}{2\pi}.
  \label{eq:rotation-count}
\end{equation}
若 $\rho(r)$ 很小，说明该维度在训练区间内尚未完整旋转，应像 PI 一样充分插值以避免 OOD；若 $\rho(r)$ 很大，说明它已经多次旋转，应尽量保留以维持局部分辨率。定义分段线性 ramp：
\begin{equation}
  \gamma(\rho)=
  \begin{cases}
    0,&\rho<\alpha,\\
    \dfrac{\rho-\alpha}{\beta-\alpha},&\alpha\le\rho\le\beta,\\
    1,&\rho>\beta,
  \end{cases}
  \label{eq:ramp}
\end{equation}
则 NTK-by-parts 的新频率可写为：
\begin{equation}
  \theta'_r=
  \bigl(1-\gamma(\rho(r))\bigr)\frac{\theta_r}{s}
  +\gamma(\rho(r))\theta_r.
  \label{eq:ntk-by-parts}
\end{equation}
低圈数频率接近 PI，高圈数频率保持原值，中间频率平滑过渡。YaRN 对 LLaMA 系列给出的经验设置为 $\alpha=1,\beta=32$，但这不是跨模型通用常数\cite{peng2024yarn}。

\subsection{YaRN：NTK-by-parts 加注意力尺度校准}

扩长后 attention entropy 往往改变。YaRN 在 NTK-by-parts 基础上加入 attention scaling\cite{peng2024yarn}：
\begin{equation}
  \alpha_{mn}=\softmax_n\!\left(
  \frac{q_m^\top k_n}{t\sqrt{d_h}}
  \right),
  \label{eq:yarn-temperature}
\end{equation}
并可通过同时缩放 $q,k$ 来等价实现：
\begin{equation}
  \tilde q=\frac{q}{\sqrt{t}},\qquad
  \tilde k=\frac{k}{\sqrt{t}}.
  \label{eq:yarn-qk-scale}
\end{equation}
论文对 LLaMA/Llama 2 给出经验关系：
\begin{equation}
  \frac{1}{\sqrt{t}}=0.1\ln s+1.
  \label{eq:yarn-fit}
\end{equation}
因此 YaRN 不是“另一种单一插值公式”，而是\term{分频率插值 $+$ 注意力温度校准}。论文报告它用比 PI 少约 $2.5\times$ 的训练步数和数据达到相近扩展效果\cite{peng2024yarn}。实际框架中的 \textnormal{factor}、\textnormal{beta\_fast}、\textnormal{beta\_slow}、\textnormal{attention\_factor} 等命名可能不同，必须与 checkpoint 训练配置一致。

\subsection{LongRoPE：逐维、分位置搜索缩放因子}

前述方法都用解析规则生成缩放。LongRoPE 进一步利用两种非均匀性\cite{ding2024longrope}：
\begin{enumerate}
  \item 不同 RoPE 维度需要不同缩放因子 $\lambda_r$；
  \item 序列开头一段位置对插值更敏感，可设置保留区间 $n_{\mathrm{hat}}$。
\end{enumerate}
其一般化频率写作：
\begin{equation}
  \theta'_r=\frac{\theta_r}{\lambda_r},
  \label{eq:longrope-scale}
\end{equation}
并用 evolutionary search 在约束空间中寻找 $\{\lambda_r\}$ 与 $n_{\mathrm{hat}}$。LongRoPE 先在 256K 长度适配，再基于该模型进行第二阶段搜索与插值，论文报告达到 2048K；同时为短上下文搜索另一套缩放配置以恢复原长度性能\cite{ding2024longrope}。关键思想是：\term{扩展倍率不应被假设为跨维度、跨位置均匀。}

\subsection{LongRoPE2：从“理论周期”转向“有效训练周期”}

LongRoPE2 指出，低频/高维 RoPE 在预训练语料中不仅理论周期长，而且真正发生长距离依赖的样本稀少，导致其\term{有效相位范围}比理论预测更窄\cite{longrope2}。因此，理论临界维度未必等于实际临界维度。它做了三项改进：
\begin{enumerate}
  \item 使用 needle-driven perplexity 指导 evolutionary search，联合寻找实际临界维度与逐维缩放因子；
  \item 对 OOD 维度施加单调非减的 $\lambda_r$ 约束；
  \item 混合上下文训练：短样本使用原始 RoPE，长样本使用 rescaled RoPE。
\end{enumerate}
混合训练目标可抽象为：
\begin{equation}
  \mathcal{L}=\eta\,\mathcal{L}_{\mathrm{short}}(\theta_{\mathrm{orig}})
  +(1-\eta)\,\mathcal{L}_{\mathrm{long}}(\theta_{\mathrm{scaled}}),
  \label{eq:mixed-context-loss}
\end{equation}
它直接处理“长上下文适配”和“短上下文保持”之间的冲突。论文在 LLaMA3-8B 与 Phi3-mini 上报告 128K 有效上下文与较高短上下文保持率；这些数字是论文实验结论，不应外推为所有模型的默认保证\cite{longrope2}。

\subsection{方法对比与适用条件}

\begin{table}[htbp]
\centering
\small
\begin{tabularx}{\linewidth}{@{}p{2.4cm}p{3.0cm}p{3.2cm}X@{}}
\toprule
方法 & 缩放对象 & 主要优点 & 主要代价/风险 \\
\midrule
PI & 所有位置/频率统一除以 $s$ & 稳定、实现最简单 & 高频局部分辨率下降，通常需微调 \\
NTK-aware & 修改 base，逐维非均匀缩放 & 保留最高频局部结构 & 社区经验起源，超参依模型 \\
Dynamic NTK & 倍率随当前长度变化 & 短输入保持原 RoPE & 与旋转后 KV Cache 冲突 \\
NTK-by-parts & 按旋转圈数分区插值 & 高频保留、低频防 OOD & 需选择 $\alpha,\beta$ \\
YaRN & NTK-by-parts + 温度校准 & 长短程折中更完整 & 必须匹配训练 recipe 与 attention factor \\
LongRoPE & 搜索逐维因子与位置保留区 & 能利用非均匀性，支持渐进扩长 & 搜索与验证成本更高 \\
LongRoPE2 & 搜索有效临界维度 + 混合训练 & 同时优化有效长度和短程保持 & 需要专门数据、搜索和 mid-training \\
\bottomrule
\end{tabularx}
\caption{主要 RoPE 扩展方法的统一比较。}
\label{tab:rope-scaling}
\end{table}

\noteBox{\textbf{实践建议：}对已有 checkpoint，应复用其原生 RoPE 配置和训练 recipe；\textnormal{rope\_theta}、PI factor、YaRN factor 与 LongRoPE 的逐维因子并不等价。从头训练长上下文模型时，需要联合设计 RoPE scaling、长短样本 mix、position-id 语义和评估矩阵。}

\section{另一种设计：ALiBi 的距离线性偏置}

ALiBi 不旋转向量，而是在第 $h$ 个 head 的注意力分数中加入线性距离惩罚\cite{press2022}：
\begin{equation}
  s_{ij}^{(h)}=
  \frac{q_i^{(h)\top}k_j^{(h)}}{\sqrt{d_h}}
  -m_h|i-j|.
  \label{eq:alibi}
\end{equation}
不同 head 使用不同斜率 $m_h$。它形式简单、没有位置表，并显式注入 recency bias；但其归纳偏置与 RoPE 的多频率相位不同，是相对位置 bias 分支而非 RoPE 扩展。

\section{三类位置机制的统一视角}

三类方法分别修改输入、token 对分数和 Query-Key 几何关系：
\begin{align}
  x_i'&=x_i+p_i, &&\text{绝对位置}, \label{eq:unified-absolute}\\
  s_{ij}&=\frac{q_i^\top k_j}{\sqrt{d_h}}+b(i-j), &&\text{相对位置 bias}, \label{eq:unified-relative}\\
  s_{ij}&=\frac{(R_iq_i)^\top(R_jk_j)}{\sqrt{d_h}}
  =\frac{q_i^\top R_{j-i}k_j}{\sqrt{d_h}}, &&\text{RoPE}. \label{eq:unified-rope}
\end{align}

\begin{table}[htbp]
\centering
\footnotesize
\begin{tabularx}{\linewidth}{@{}p{2.2cm}p{1.5cm}p{1.8cm}p{1.8cm}X@{}}
\toprule
机制 & 注入位置 & 相对距离 & 超长位置可计算 & 可靠外推条件 \\
\midrule
Learned Absolute & 输入 & 不显式 & 通常否 & 受位置表与训练长度限制 \\
Sinusoidal & 输入 & 几何中隐含 & 是 & 可计算不等于可利用 \\
Relative Bias & logits/K/V & 显式 & 依分桶而定 & 取决于分桶与训练 \\
RoPE & Q/K & 点积中显式 & 是 & 通常需 scaling 与适配训练 \\
ALiBi & logits & 显式 & 是 & 具有强 recency bias \\
\bottomrule
\end{tabularx}
\caption{位置机制的结构比较。}
\label{tab:position-compare}
\end{table}

\section{工程实验与评估建议}

\subsection{配置必须随 checkpoint 固化}

对 decoder-only LLM，至少应保存并验证：
\begin{itemize}
  \item \textnormal{rope\_theta} 或 base frequency；
  \item rotary dimension：全维或 partial rotary；
  \item 维度配对布局；
  \item 原始训练长度与目标长度；
  \item scaling 类型、factor、attention factor 与逐维参数；
  \item packing、prefix cache、sliding window 下 position id 的定义。
\end{itemize}

\subsection{长上下文评估不能只看“能否输入”}

建议至少报告：
\begin{enumerate}
  \item 短上下文保持率：扩长前后原长度 PPL 与 benchmark；
  \item 按 token position、依赖跨度分桶的 PPL；
  \item needle/passkey 的长度--深度二维网格；
  \item RULER 类多任务长上下文评估，而非单一检索；
  \item 有效上下文：性能开始显著退化的位置，而非配置最大长度；
  \item prefill/decode latency、峰值显存、KV Cache 占用与吞吐。
\end{enumerate}

代码模型还应补充 definition-use 跨度、跨文件 import/call graph 定位、同名符号消歧、repo-level completion 与 patch correctness，并验证扩长后 HumanEval 类短程语法和算法能力是否退化。

\section{常见误解}

\begin{enumerate}
  \item \textbf{Self-Attention 完全看不见顺序。}更准确地说，不含位置机制的 Self-Attention 对输入置换等变。causal mask 提供可见方向，但不等于精细距离表示。
  \item \textbf{正弦编码是绝对位置，所以不含相对信息。}式~\eqref{eq:sin-shift} 与式~\eqref{eq:sin-dot} 表明固定偏移对应固定旋转，点积只依赖 $m-n$。
  \item \textbf{RoPE 是相对位置编码，因此不编码绝对位置。}单个 $R_mq_m$ 明确依赖 $m$；只有点积时绝对相位才相消。
  \item \textbf{位置函数能算到 128K，模型就有 128K 上下文。}配置长度、数值可运行长度与有效上下文长度是三件事。
  \item \textbf{RoPE 天然让注意力随距离单调下降。}式~\eqref{eq:rope-oscillation} 是周期函数，不存在逐样本严格单调保证。
  \item \textbf{只调大 \textnormal{rope\_theta} 就完成扩长。}频率缩放只是第一步，还要验证模型适配、KV Cache 一致性、短程保持与有效上下文。
\end{enumerate}

\section{总结与展望}

绝对位置编码把坐标写入 token 表示，相对位置编码直接修改 token 对之间的关系，RoPE 则通过 Query-Key 旋转把绝对位置转化为点积中的相对相位。PI、NTK-aware、YaRN 和 LongRoPE 进一步处理频率缩放与长短上下文兼容问题；RoPE scaling law 则从周期覆盖出发，联系 base、训练长度、临界维度与外推范围。不过，现有位置机制仍有若干问题没有解决：

\begin{enumerate}
  \item \textbf{标称长度与有效长度不一致。}能接收 128K 或 1M token，不代表能稳定利用所有位置；有效上下文还受训练依赖跨度、优化过程和注意力稀释影响。
  \item \textbf{周期性带来相位混叠。}不同距离可能产生相近相位；多频率组合能够缓解，却不能从理论上消除这一问题。
  \item \textbf{扩长与短程能力仍有冲突。}插值会降低局部分辨率，逐维缩放和混合长度训练也缺少可跨模型迁移的稳定规律。
  \item \textbf{系统实现可能破坏位置语义。}packing、prefix cache、sliding window、speculative decoding 和 KV Cache 复用都可能造成隐蔽的 position-id 错位。
  \item \textbf{一维距离难以表达复杂结构。}代码的文件层级、调用图和 definition-use 关系，以及多模态数据的空间与时间结构，都超出单一序列坐标的表达范围。
  \item \textbf{位置编码与高效注意力缺少统一设计。}稀疏注意力、滑窗、KV 压缩和线性注意力会改变信息传播路径，位置机制需要与算子共同设计。
  \item \textbf{评估方法仍然偏弱。}needle-in-a-haystack 不能充分反映多跳推理、跨文件代码理解和长程因果依赖。
\end{enumerate}

后续研究不应只放大标称上下文，而应建立可预测、可训练、可缓存且能表达多维结构的位置机制，并用真实长依赖任务检验远距离信息利用率。对代码模型而言，definition-use、跨文件调用和 repository-level 修改比随机 needle 更接近实际需求。

\clearpage
\appendix
\section{代表性 LLM 的位置编码选择}

表~\ref{tab:llm-position-encoding} 汇总公开论文、技术报告或官方开源配置中可以核验的位置机制，资料状态更新至 2026 年 8 月 7 日。表中的“上下文扩展”只记录位置编码相关做法，不代表模型仅凭该机制就能达到标称上下文长度。对于没有公开架构细节的模型，不能依据上下文长度或同系列旧模型反推出其位置编码。

\begin{longtable}{@{}p{2.7cm}p{3.6cm}p{8.0cm}@{}}
\caption{代表性语言模型公开披露的位置编码方式。}
\label{tab:llm-position-encoding}\\
\toprule
模型/系列 & 位置机制 & 技术报告中的相关设计 \\
\midrule
\endfirsthead
\multicolumn{3}{l}{\small 表~\thetable（续）}\\
\toprule
模型/系列 & 位置机制 & 技术报告中的相关设计 \\
\midrule
\endhead
\midrule
\multicolumn{3}{r}{\small 下页续}\\
\endfoot
\bottomrule
\endlastfoot

GPT-2 / GPT-3 & 可学习绝对位置 embedding & GPT-2 使用可学习的位置 embedding；GPT-3 沿用 GPT-2 的基本 Transformer 架构，因此仍采用 learned absolute 方案\cite{radford2019gpt2,brown2020gpt3}。这类设计的最大长度通常受位置表和训练长度约束。\\

T5 & 分桶相对位置 bias & 每个 attention head 对相对距离桶学习标量偏置；近距离细分、远距离近似对数分桶\cite{raffel2020}。\\

BLOOM & ALiBi & 不向 embedding 加位置向量，而是在注意力 logits 上施加按距离增长的线性负偏置\cite{lescao2022bloom}。\\

Falcon & ALiBi & Falcon 系列技术报告采用 ALiBi；其消融实验同时比较了 ALiBi、URPE 与 RoPE\cite{almazrouei2023falcon}。\\

PaLM & RoPE & 使用 RoPE 取代绝对或相对位置 embedding；报告将其长序列表现列为采用理由之一\cite{chowdhery2022palm}。\\

LLaMA / Llama 2 & RoPE & LLaMA 以 RoPE 取代绝对位置 embedding；Llama 2 延续 RoPE，并把预训练上下文从 2K 扩为 4K\cite{touvron2023llama,touvron2023llama2}。\\

Code Llama & RoPE，增大基数 & 在 Llama 2 基础上把 RoPE base 调为 $10^6$，使用 16K 序列继续训练；报告研究了向 100K 的外推\cite{roziere2023codellama}。\\

Llama 3 & RoPE，$\theta=500{,}000$ & 技术报告把 RoPE base 提高到 500K，用于改善长上下文支持；该配置对 8B、70B 与 405B 模型一致\cite{dubey2024llama3}。\\

Qwen2 & RoPE + YaRN + DCA & 预训练后期把上下文从 4K 扩到 32K，并把 RoPE base 从 10K 调为 $10^6$；推理侧结合 YaRN 与 Dual Chunk Attention 支持更长输入\cite{yang2024qwen2}。\\

Qwen3 / Qwen3.6 & RoPE 系列 + YaRN & Qwen3 的官方配置使用 $\theta=10^6$，可通过 YaRN 从原生 32K 扩展到 131K\cite{yang2025qwen3}。Qwen3.6 进一步采用全注意力与线性注意力混合结构；其官方配置在全注意力层使用部分旋转的 M-RoPE，$\theta=10^7$，原生支持 262K，并给出使用 YaRN 扩展至约 1M 的配置\cite{qwen36config}。\\

Gemma & RoPE & Gemma 在每层使用 rotary positional embeddings，而不是绝对位置 embedding；基础模型训练长度为 8K\cite{gemmateam2024}。\\

Phi-3 / Phi-3.5 & RoPE 系架构 + LongRoPE & Phi-3-mini 采用与 Llama 2 相近的 block；128K 版本通过 LongRoPE 扩展，Phi-3.5 继续使用 LongRoPE 与混合上下文训练\cite{abdin2024phi3}。\\

DeepSeek-V2 / V3 & MLA 中的 decoupled RoPE & 为兼容低秩 KV 压缩，将携带 RoPE 的 Query/Key 子空间与内容压缩子空间解耦；V3 延续该 MLA 设计\cite{deepseek2024v2,deepseek2024v3}。\\

LongCat 系列 & MLA 中的 decoupled RoPE & LongCat-Flash、LongCat-Flash-Thinking 与 2601 版本沿用 MLA，把 Query/Key 分成 RoPE 与 NoPE 子空间。初版 Flash 的官方配置为 64 维 RoPE 子空间、128 维 NoPE 子空间，$\theta=10^7$，最大位置为 131K；2601 的公开配置将 $\theta$ 设为 $10^6$\cite{longcatflash,longcatthinking2601,longcatconfigs}。\\

Kimi K3 & NoPE + Kimi Delta Attention & 不向 MLA 的 Query/Key 注入显式位置编码；位置敏感性由 KDA 的递归门控与衰减机制隐式提供。技术报告指出该设计无需 RoPE rescaling 或 interpolation 即可继续训练到 1M 上下文\cite{kimik3}。\\

GLM-5 & MLA/DSA 中的 decoupled RoPE & GLM-5 以 DSA 稀疏化 MLA 的注意力计算；官方配置保留 64 维 RoPE 子空间与 192 维 NoPE 子空间，并设置 $\theta=10^6$。DSA 改变稀疏检索路径，但没有取消 RoPE\cite{glm5,glm5config}。\\

Qwen3.8-Max & 未公开 & 官方服务文档已列出 Qwen3.8-Max，但截至本文资料截止日尚无公开技术报告、权重或配置可以确认位置机制。因此不能直接套用 Qwen3/Qwen3.6 的 RoPE 方案\cite{qwen38docs}。\\

GPT-4 & 未公开 & GPT-4 技术报告未披露模型结构、位置编码或 RoPE 配置，因此无法从公开材料确定其位置机制\cite{openai2023gpt4}。\\

\end{longtable}

从公开资料可以看到，decoder-only 开源模型在 2022 年以后明显集中到 RoPE，但并未完全收敛：BLOOM 与 Falcon 使用 ALiBi，T5 系模型采用相对位置 bias，DeepSeek、LongCat 与 GLM-5 为适配 MLA 或稀疏注意力而拆分 RoPE/NoPE 子空间，Kimi K3 则转向由 KDA 隐式提供位置信息的 NoPE 方案。所谓“使用 RoPE”也不是一个完整配置，base、rotary dimension、scaling、训练长度和 position-id 语义都会改变实际行为；对于 Qwen3.8-Max 这类未披露架构的模型，表中保留“未知”比沿用系列历史配置更可靠。

\clearpage
\begin{thebibliography}{99}
\small
\setlength{\itemsep}{1pt}

\bibitem{vaswani2017}
A. Vaswani et al. \href{https://arxiv.org/abs/1706.03762}{\textit{Attention Is All You Need}}. NeurIPS, 2017.

\bibitem{devlin2019}
J. Devlin et al. \href{https://arxiv.org/abs/1810.04805}{\textit{BERT: Pre-training of Deep Bidirectional Transformers for Language Understanding}}. NAACL, 2019.

\bibitem{shaw2018}
P. Shaw, J. Uszkoreit, A. Vaswani. \href{https://aclanthology.org/N18-2074/}{\textit{Self-Attention with Relative Position Representations}}. NAACL, 2018.

\bibitem{dai2019}
Z. Dai et al. \href{https://arxiv.org/abs/1901.02860}{\textit{Transformer-XL: Attentive Language Models Beyond a Fixed-Length Context}}. ACL, 2019.

\bibitem{raffel2020}
C. Raffel et al. \href{https://www.jmlr.org/papers/v21/20-074.html}{\textit{Exploring the Limits of Transfer Learning with a Unified Text-to-Text Transformer}}. JMLR, 2020.

\bibitem{su2021}
J. Su et al. \href{https://arxiv.org/abs/2104.09864}{\textit{RoFormer: Enhanced Transformer with Rotary Position Embedding}}. 2021.

\bibitem{press2022}
O. Press, N. A. Smith, M. Lewis. \href{https://arxiv.org/abs/2108.12409}{\textit{Train Short, Test Long: Attention with Linear Biases Enables Input Length Extrapolation}}. ICLR, 2022.

\bibitem{chen2023pi}
S. Chen et al. \href{https://arxiv.org/abs/2306.15595}{\textit{Extending Context Window of Large Language Models via Position Interpolation}}. 2023.

\bibitem{liu2024scalingrope}
X. Liu et al. \href{https://arxiv.org/abs/2310.05209}{\textit{Scaling Laws of RoPE-based Extrapolation}}. ICLR, 2024.

\bibitem{peng2024yarn}
B. Peng et al. \href{https://arxiv.org/abs/2309.00071}{\textit{YaRN: Efficient Context Window Extension of Large Language Models}}. ICLR, 2024; arXiv revised 2026.

\bibitem{ding2024longrope}
Y. Ding et al. \href{https://arxiv.org/abs/2402.13753}{\textit{LongRoPE: Extending LLM Context Window Beyond 2 Million Tokens}}. ICML, 2024.

\bibitem{longrope2}
N. Shang et al. \href{https://arxiv.org/abs/2502.20082}{\textit{LongRoPE2: Near-Lossless LLM Context Window Scaling}}. 2025.

\bibitem{radford2019gpt2}
A. Radford et al. \href{https://cdn.openai.com/better-language-models/language_models_are_unsupervised_multitask_learners.pdf}{\textit{Language Models are Unsupervised Multitask Learners}}. OpenAI, 2019.

\bibitem{brown2020gpt3}
T. Brown et al. \href{https://arxiv.org/abs/2005.14165}{\textit{Language Models are Few-Shot Learners}}. NeurIPS, 2020.

\bibitem{lescao2022bloom}
T. Le Scao et al. \href{https://arxiv.org/abs/2211.05100}{\textit{BLOOM: A 176B-Parameter Open-Access Multilingual Language Model}}. 2022.

\bibitem{almazrouei2023falcon}
E. Almazrouei et al. \href{https://arxiv.org/abs/2311.16867}{\textit{The Falcon Series of Open Language Models}}. 2023.

\bibitem{chowdhery2022palm}
A. Chowdhery et al. \href{https://arxiv.org/abs/2204.02311}{\textit{PaLM: Scaling Language Modeling with Pathways}}. 2022.

\bibitem{touvron2023llama}
H. Touvron et al. \href{https://arxiv.org/abs/2302.13971}{\textit{LLaMA: Open and Efficient Foundation Language Models}}. 2023.

\bibitem{touvron2023llama2}
H. Touvron et al. \href{https://arxiv.org/abs/2307.09288}{\textit{Llama 2: Open Foundation and Fine-Tuned Chat Models}}. 2023.

\bibitem{roziere2023codellama}
B. Rozière et al. \href{https://arxiv.org/abs/2308.12950}{\textit{Code Llama: Open Foundation Models for Code}}. 2023.

\bibitem{dubey2024llama3}
A. Dubey et al. \href{https://arxiv.org/abs/2407.21783}{\textit{The Llama 3 Herd of Models}}. 2024.

\bibitem{yang2024qwen2}
A. Yang et al. \href{https://arxiv.org/abs/2407.10671}{\textit{Qwen2 Technical Report}}. 2024.

\bibitem{gemmateam2024}
Gemma Team. \href{https://arxiv.org/abs/2403.08295}{\textit{Gemma: Open Models Based on Gemini Research and Technology}}. 2024.

\bibitem{abdin2024phi3}
M. Abdin et al. \href{https://arxiv.org/abs/2404.14219}{\textit{Phi-3 Technical Report: A Highly Capable Language Model Locally on Your Phone}}. 2024.

\bibitem{deepseek2024v2}
DeepSeek-AI. \href{https://arxiv.org/abs/2405.04434}{\textit{DeepSeek-V2: A Strong, Economical, and Efficient Mixture-of-Experts Language Model}}. 2024.

\bibitem{deepseek2024v3}
DeepSeek-AI. \href{https://arxiv.org/abs/2412.19437}{\textit{DeepSeek-V3 Technical Report}}. 2024.

\bibitem{yang2025qwen3}
Qwen Team. \href{https://arxiv.org/abs/2505.09388}{\textit{Qwen3 Technical Report}}. 2025.

\bibitem{qwen36config}
Qwen Team. \href{https://huggingface.co/Qwen/Qwen3.6-35B-A3B-FP8/blob/main/config.json}{\textit{Qwen3.6-35B-A3B Official Model Configuration}}; \href{https://huggingface.co/Qwen/Qwen3.6-35B-A3B/blob/main/README.md}{\textit{Processing Ultra-Long Texts}}. 2026.

\bibitem{longcatflash}
Meituan LongCat Team. \href{https://arxiv.org/abs/2509.01322}{\textit{LongCat-Flash Technical Report}}. 2025.

\bibitem{longcatthinking2601}
Meituan LongCat Team. \href{https://arxiv.org/abs/2601.16725}{\textit{LongCat-Flash-Thinking-2601 Technical Report}}. 2026.

\bibitem{longcatconfigs}
Meituan LongCat Team. \href{https://huggingface.co/meituan-longcat/LongCat-Flash-Chat/blob/main/config.json}{\textit{LongCat-Flash-Chat Official Model Configuration}}; \href{https://huggingface.co/meituan-longcat/LongCat-Flash-Thinking-2601-FP8/blob/main/config.json}{\textit{LongCat-Flash-Thinking-2601 Official Model Configuration}}. 2025--2026.

\bibitem{kimik3}
Kimi Team. \href{https://arxiv.org/abs/2607.24653}{\textit{Kimi K3: Open Frontier Intelligence}}. 2026.

\bibitem{glm5}
GLM-5 Team. \href{https://arxiv.org/abs/2602.15763}{\textit{GLM-5: From Vibe Coding to Agentic Engineering}}. 2026.

\bibitem{glm5config}
Z.ai. \href{https://huggingface.co/zai-org/GLM-5/blob/main/config.json}{\textit{GLM-5 Official Model Configuration}}. 2026.

\bibitem{qwen38docs}
Alibaba Cloud. \href{https://help.aliyun.com/en/model-studio/models}{\textit{Model Studio: Supported Models}}. Accessed 2026-08-07.

\bibitem{openai2023gpt4}
OpenAI. \href{https://arxiv.org/abs/2303.08774}{\textit{GPT-4 Technical Report}}. 2023.

\end{thebibliography}

\end{document}
